% 07_appendix.tex —— 附录
\section{程序代码}

本文的完整求解、验证和绘图脚本保存在工程目录中。本附录列出关键代码片段，完整脚本路径如下：
\begin{itemize}
  \item \path{baseline_v1/code/model_solver.py}：MILP 建模求解并导出四问调度表；
  \item \path{code/analysis/verify_optimality_chain.py}：Q2 下界与 Q4 采购模型验证；
  \item \path{code/plot/plot_gantt.py}：根据 CSV 结果表重新生成论文图件。
\end{itemize}

\lstset{basicstyle=\scriptsize\ttfamily, breakatwhitespace=false}

\subsection{MILP 调度模型核心代码}
\lstinputlisting[
  language=Python,
  firstline=165,
  lastline=284,
  caption={MILP 变量、目标与约束构造核心代码}
]{../../baseline_v1/code/model_solver.py}

\subsection{结果导出与四问求解入口}
\lstinputlisting[
  language=Python,
  firstline=326,
  lastline=386,
  caption={四问求解流程与结果导出代码}
]{../../baseline_v1/code/model_solver.py}

\subsection{CP-SAT 最优性验证核心代码}
\lstinputlisting[
  language=Python,
  firstline=69,
  lastline=134,
  caption={问题二双瓶颈资源松弛下界模型}
]{../../code/analysis/verify_optimality_chain.py}

\lstinputlisting[
  language=Python,
  firstline=269,
  lastline=318,
  caption={问题四显式采购模型的两阶段优化代码}
]{../../code/analysis/verify_optimality_chain.py}

\subsection{论文图件生成代码}
\lstinputlisting[
  language=Python,
  firstline=53,
  lastline=145,
  caption={甘特图与工期汇总图生成代码}
]{../../code/plot/plot_gantt.py}

\section{补充图表}

本节给出正文未完整展开的题目要求结果表。表格均由 \path{baseline_v1/results/} 中的 CSV 文件自动生成：问题二、三、四的完整设备调度表分别对应 \path{table2_question2.csv}、\path{table3_question3.csv}、\path{table4_question4.csv}，设备购置表对应 \path{table5_purchase_plan.csv}，四问汇总表对应 \path{summary.csv}。

% This file is generated from baseline_v1/results/*.csv. Do not edit numbers by hand.

\begingroup
\scriptsize
\renewcommand{\arraystretch}{1.08}
\begin{longtable}{@{}r p{3.6cm} c c r c@{}}
\caption{问题二完整设备调度表}\label{tab:appendix-q2}\\
\toprule
序号 & 设备编号 & 起始时间 & 结束时间 & 持续工作时间(s) & 工序编号 \\
\midrule
\endfirsthead
\multicolumn{6}{c}{续表：问题二完整设备调度表}\\
\toprule
序号 & 设备编号 & 起始时间 & 结束时间 & 持续工作时间(s) & 工序编号 \\
\midrule
\endhead
\bottomrule
\endfoot
1 & 精密灌装机1-4 & 00:03:20 & 01:33:20 & 5400 & A1 \\
2 & 自动化输送臂1-2 & 00:03:20 & 01:15:20 & 4320 & A1 \\
3 & 工业清洗机1-4 & 00:03:20 & 04:03:20 & 14400 & E1 \\
4 & 工业清洗机1-2 & 00:03:50 & 02:56:38 & 10368 & C1 \\
5 & 自动化输送臂1-3 & 00:03:50 & 02:56:38 & 10368 & C1 \\
6 & 工业清洗机1-5 & 00:05:10 & 01:17:10 & 4320 & B1 \\
7 & 工业清洗机1-3 & 00:05:55 & 02:29:55 & 8640 & D1 \\
8 & 精密灌装机1-2 & 01:17:10 & 08:47:10 & 27000 & B2 \\
9 & 自动化输送臂1-1 & 01:17:10 & 06:17:10 & 18000 & B2 \\
10 & 工业清洗机1-1 & 01:33:20 & 03:33:20 & 7200 & A2 \\
11 & 高速抛光机1-1 & 01:33:20 & 06:33:20 & 18000 & A2 \\
12 & 精密灌装机1-4 & 02:29:55 & 06:29:55 & 14400 & D2 \\
13 & 自动化输送臂1-4 & 02:29:55 & 05:09:55 & 9600 & D2 \\
14 & 精密灌装机1-5 & 02:56:38 & 05:00:04 & 7406 & C2 \\
15 & 精密灌装机1-1 & 04:03:20 & 05:46:12 & 6172 & E2 \\
16 & 精密灌装机1-3 & 05:00:04 & 06:48:04 & 6480 & C3\_R1 \\
17 & 自动化输送臂1-2 & 05:00:04 & 06:26:28 & 5184 & C3\_R1 \\
18 & 工业清洗机1-2 & 05:46:12 & 11:46:12 & 21600 & E3 \\
19 & 自动传感多功能机1-1 & 05:46:12 & 07:46:12 & 7200 & E3 \\
20 & 精密灌装机1-5 & 06:29:55 & 07:47:04 & 4629 & D3 \\
21 & 工业清洗机1-4 & 06:48:04 & 10:48:04 & 14400 & C4\_R1 \\
22 & 高速抛光机1-1 & 06:48:04 & 10:08:04 & 12000 & C4\_R1 \\
23 & 精密灌装机1-5 & 08:47:10 & 09:48:53 & 3703 & B3 \\
24 & 自动传感多功能机1-1 & 10:12:24 & 15:12:24 & 18000 & D4 \\
25 & 高速抛光机1-1 & 10:12:24 & 22:42:24 & 45000 & D4 \\
26 & 自动传感多功能机1-1 & 15:16:44 & 19:16:44 & 14400 & C5\_R1 \\
27 & 精密灌装机1-1 & 19:16:44 & 21:04:44 & 6480 & C3\_R2 \\
28 & 自动化输送臂1-1 & 19:16:44 & 20:43:08 & 5184 & C3\_R2 \\
29 & 自动传感多功能机1-1 & 22:42:24 & 27:42:24 & 18000 & D5 \\
30 & 工业清洗机1-3 & 22:46:44 & 26:46:44 & 14400 & C4\_R2 \\
31 & 高速抛光机1-1 & 22:46:44 & 26:06:44 & 12000 & C4\_R2 \\
32 & 自动传感多功能机1-1 & 27:46:44 & 31:46:44 & 14400 & C5\_R2 \\
33 & 精密灌装机1-1 & 31:46:44 & 33:34:44 & 6480 & C3\_R3 \\
34 & 自动化输送臂1-4 & 31:46:44 & 33:13:08 & 5184 & C3\_R3 \\
35 & 自动传感多功能机1-1 & 31:55:54 & 35:31:54 & 12960 & B4 \\
36 & 高速抛光机1-1 & 31:55:54 & 34:55:54 & 10800 & B4 \\
37 & 工业清洗机1-5 & 35:05:04 & 39:05:04 & 14400 & C4\_R3 \\
38 & 高速抛光机1-1 & 35:05:04 & 38:25:04 & 12000 & C4\_R3 \\
39 & 自动传感多功能机1-1 & 35:40:24 & 40:40:24 & 18000 & A3 \\
40 & 高速抛光机1-1 & 38:29:24 & 45:29:24 & 25200 & D6 \\
41 & 自动传感多功能机1-1 & 40:49:09 & 44:49:09 & 14400 & C5\_R3 \\
\end{longtable}
\endgroup

\begingroup
\scriptsize
\renewcommand{\arraystretch}{1.08}
\begin{longtable}{@{}r p{3.4cm} c c r c c@{}}
\caption{问题三完整设备调度表}\label{tab:appendix-q3}\\
\toprule
序号 & 设备编号 & 起始时间 & 结束时间 & 持续工作时间(s) & 工序编号 & 班组 \\
\midrule
\endfirsthead
\multicolumn{7}{c}{续表：问题三完整设备调度表}\\
\toprule
序号 & 设备编号 & 起始时间 & 结束时间 & 持续工作时间(s) & 工序编号 & 班组 \\
\midrule
\endhead
\bottomrule
\endfoot
1 & 精密灌装机1-1 & 00:03:20 & 01:33:20 & 5400 & A1 & 1 \\
2 & 自动化输送臂1-1 & 00:03:20 & 01:15:20 & 4320 & A1 & 1 \\
3 & 工业清洗机1-4 & 00:03:20 & 04:03:20 & 14400 & E1 & 1 \\
4 & 工业清洗机2-2 & 00:03:50 & 01:15:50 & 4320 & B1 & 2 \\
5 & 工业清洗机1-5 & 00:03:50 & 02:56:38 & 10368 & C1 & 1 \\
6 & 自动化输送臂1-2 & 00:03:50 & 02:56:38 & 10368 & C1 & 1 \\
7 & 工业清洗机2-4 & 00:05:40 & 02:29:40 & 8640 & D1 & 2 \\
8 & 精密灌装机2-2 & 01:15:50 & 08:45:50 & 27000 & B2 & 2 \\
9 & 自动化输送臂1-4 & 01:15:50 & 06:15:50 & 18000 & B2 & 1 \\
10 & 工业清洗机2-2 & 01:33:20 & 03:33:20 & 7200 & A2 & 2 \\
11 & 高速抛光机2-1 & 01:33:20 & 06:33:20 & 18000 & A2 & 2 \\
12 & 精密灌装机2-1 & 02:29:40 & 06:29:40 & 14400 & D2 & 2 \\
13 & 自动化输送臂2-2 & 02:29:40 & 05:09:40 & 9600 & D2 & 2 \\
14 & 精密灌装机2-5 & 02:56:38 & 05:00:04 & 7406 & C2 & 2 \\
15 & 精密灌装机1-3 & 04:03:20 & 05:46:12 & 6172 & E2 & 1 \\
16 & 精密灌装机1-5 & 05:00:04 & 06:48:04 & 6480 & C3\_R1 & 1 \\
17 & 自动化输送臂1-2 & 05:00:04 & 06:26:28 & 5184 & C3\_R1 & 1 \\
18 & 工业清洗机1-2 & 05:46:12 & 11:46:12 & 21600 & E3 & 1 \\
19 & 自动传感多功能机1-1 & 05:46:12 & 07:46:12 & 7200 & E3 & 1 \\
20 & 精密灌装机1-2 & 06:29:40 & 07:46:49 & 4629 & D3 & 1 \\
21 & 工业清洗机2-3 & 06:48:04 & 10:48:04 & 14400 & C4\_R1 & 2 \\
22 & 高速抛光机1-1 & 06:48:04 & 10:08:04 & 12000 & C4\_R1 & 1 \\
23 & 自动传感多功能机2-1 & 07:46:49 & 12:46:49 & 18000 & D4 & 2 \\
24 & 高速抛光机2-1 & 07:46:49 & 20:16:49 & 45000 & D4 & 2 \\
25 & 精密灌装机2-5 & 08:45:50 & 09:47:33 & 3703 & B3 & 2 \\
26 & 自动传感多功能机1-1 & 10:48:04 & 14:48:04 & 14400 & C5\_R1 & 1 \\
27 & 自动传感多功能机2-1 & 13:00:24 & 16:36:24 & 12960 & B4 & 2 \\
28 & 高速抛光机1-1 & 13:00:24 & 16:00:24 & 10800 & B4 & 1 \\
29 & 精密灌装机1-4 & 14:48:04 & 16:36:04 & 6480 & C3\_R2 & 1 \\
30 & 自动化输送臂2-3 & 14:48:04 & 16:14:28 & 5184 & C3\_R2 & 2 \\
31 & 自动传感多功能机1-1 & 14:56:49 & 19:56:49 & 18000 & A3 & 1 \\
32 & 工业清洗机1-3 & 16:36:04 & 20:36:04 & 14400 & C4\_R2 & 1 \\
33 & 高速抛光机1-1 & 16:36:04 & 19:56:04 & 12000 & C4\_R2 & 1 \\
34 & 自动传感多功能机2-1 & 20:16:49 & 25:16:49 & 18000 & D5 & 2 \\
35 & 自动传感多功能机1-1 & 20:36:04 & 24:36:04 & 14400 & C5\_R2 & 1 \\
36 & 精密灌装机1-5 & 24:36:04 & 26:24:04 & 6480 & C3\_R3 & 1 \\
37 & 自动化输送臂1-3 & 24:36:04 & 26:02:28 & 5184 & C3\_R3 & 1 \\
38 & 高速抛光机2-1 & 25:16:49 & 32:16:49 & 25200 & D6 & 2 \\
39 & 工业清洗机2-2 & 26:24:04 & 30:24:04 & 14400 & C4\_R3 & 2 \\
40 & 高速抛光机1-1 & 26:24:04 & 29:44:04 & 12000 & C4\_R3 & 1 \\
41 & 自动传感多功能机1-1 & 30:24:04 & 34:24:04 & 14400 & C5\_R3 & 1 \\
\end{longtable}
\endgroup

\begingroup
\scriptsize
\renewcommand{\arraystretch}{1.08}
\begin{longtable}{@{}r p{3.4cm} c c r c c@{}}
\caption{问题四完整设备调度表}\label{tab:appendix-q4}\\
\toprule
序号 & 设备编号 & 起始时间 & 结束时间 & 持续工作时间(s) & 工序编号 & 班组 \\
\midrule
\endfirsthead
\multicolumn{7}{c}{续表：问题四完整设备调度表}\\
\toprule
序号 & 设备编号 & 起始时间 & 结束时间 & 持续工作时间(s) & 工序编号 & 班组 \\
\midrule
\endhead
\bottomrule
\endfoot
1 & 精密灌装机1-1 & 00:03:20 & 01:33:20 & 5400 & A1 & 1 \\
2 & 自动化输送臂1-1 & 00:03:20 & 01:15:20 & 4320 & A1 & 1 \\
3 & 工业清洗机1-4 & 00:03:20 & 04:03:20 & 14400 & E1 & 1 \\
4 & 工业清洗机2-2 & 00:03:50 & 01:15:50 & 4320 & B1 & 2 \\
5 & 工业清洗机1-5 & 00:03:50 & 02:56:38 & 10368 & C1 & 1 \\
6 & 自动化输送臂1-2 & 00:03:50 & 02:56:38 & 10368 & C1 & 1 \\
7 & 工业清洗机2-4 & 00:05:40 & 02:29:40 & 8640 & D1 & 2 \\
8 & 精密灌装机2-2 & 01:15:50 & 08:45:50 & 27000 & B2 & 2 \\
9 & 自动化输送臂1-4 & 01:15:50 & 06:15:50 & 18000 & B2 & 1 \\
10 & 工业清洗机2-2 & 01:33:20 & 03:33:20 & 7200 & A2 & 2 \\
11 & 高速抛光机2-1 & 01:33:20 & 06:33:20 & 18000 & A2 & 2 \\
12 & 精密灌装机2-1 & 02:29:40 & 06:29:40 & 14400 & D2 & 2 \\
13 & 自动化输送臂2-2 & 02:29:40 & 05:09:40 & 9600 & D2 & 2 \\
14 & 精密灌装机2-5 & 02:56:38 & 05:00:04 & 7406 & C2 & 2 \\
15 & 精密灌装机1-3 & 04:03:20 & 05:46:12 & 6172 & E2 & 1 \\
16 & 精密灌装机1-5 & 05:00:04 & 06:48:04 & 6480 & C3\_R1 & 1 \\
17 & 自动化输送臂1-2 & 05:00:04 & 06:26:28 & 5184 & C3\_R1 & 1 \\
18 & 工业清洗机1-2 & 05:46:12 & 11:46:12 & 21600 & E3 & 1 \\
19 & 自动传感多功能机1-1 & 05:46:12 & 07:46:12 & 7200 & E3 & 1 \\
20 & 精密灌装机1-2 & 06:29:40 & 07:46:49 & 4629 & D3 & 1 \\
21 & 工业清洗机2-3 & 06:48:04 & 10:48:04 & 14400 & C4\_R1 & 2 \\
22 & 高速抛光机1-1 & 06:48:04 & 10:08:04 & 12000 & C4\_R1 & 1 \\
23 & 自动传感多功能机2-1 & 07:46:49 & 12:46:49 & 18000 & D4 & 2 \\
24 & 高速抛光机2-1 & 07:46:49 & 20:16:49 & 45000 & D4 & 2 \\
25 & 精密灌装机2-5 & 08:45:50 & 09:47:33 & 3703 & B3 & 2 \\
26 & 自动传感多功能机1-1 & 10:48:04 & 14:48:04 & 14400 & C5\_R1 & 1 \\
27 & 自动传感多功能机2-1 & 13:00:24 & 16:36:24 & 12960 & B4 & 2 \\
28 & 高速抛光机1-1 & 13:00:24 & 16:00:24 & 10800 & B4 & 1 \\
29 & 精密灌装机1-4 & 14:48:04 & 16:36:04 & 6480 & C3\_R2 & 1 \\
30 & 自动化输送臂2-3 & 14:48:04 & 16:14:28 & 5184 & C3\_R2 & 2 \\
31 & 自动传感多功能机1-1 & 14:56:49 & 19:56:49 & 18000 & A3 & 1 \\
32 & 工业清洗机1-3 & 16:36:04 & 20:36:04 & 14400 & C4\_R2 & 1 \\
33 & 高速抛光机1-1 & 16:36:04 & 19:56:04 & 12000 & C4\_R2 & 1 \\
34 & 自动传感多功能机2-1 & 20:16:49 & 25:16:49 & 18000 & D5 & 2 \\
35 & 自动传感多功能机1-1 & 20:36:04 & 24:36:04 & 14400 & C5\_R2 & 1 \\
36 & 精密灌装机1-5 & 24:36:04 & 26:24:04 & 6480 & C3\_R3 & 1 \\
37 & 自动化输送臂1-3 & 24:36:04 & 26:02:28 & 5184 & C3\_R3 & 1 \\
38 & 高速抛光机2-1 & 25:16:49 & 32:16:49 & 25200 & D6 & 2 \\
39 & 工业清洗机2-2 & 26:24:04 & 30:24:04 & 14400 & C4\_R3 & 2 \\
40 & 高速抛光机1-1 & 26:24:04 & 29:44:04 & 12000 & C4\_R3 & 1 \\
41 & 自动传感多功能机1-1 & 30:24:04 & 34:24:04 & 14400 & C5\_R3 & 1 \\
\end{longtable}
\endgroup

\begingroup
\scriptsize
\renewcommand{\arraystretch}{1.08}
\begin{longtable}{@{}p{4.0cm} r r r r@{}}
\caption{问题四设备购买表}\label{tab:appendix-purchase}\\
\toprule
设备名称 & 班组1购买台数 & 班组2购买台数 & 设备单价(元) & 小计(元) \\
\midrule
\endfirsthead
\multicolumn{5}{c}{续表：问题四设备购买表}\\
\toprule
设备名称 & 班组1购买台数 & 班组2购买台数 & 设备单价(元) & 小计(元) \\
\midrule
\endhead
\bottomrule
\endfoot
自动化输送臂 & 0 & 0 & 50000 & 0 \\
工业清洗机 & 0 & 0 & 40000 & 0 \\
精密灌装机 & 0 & 0 & 35000 & 0 \\
自动传感多功能机 & 0 & 0 & 80000 & 0 \\
高速抛光机 & 0 & 0 & 75000 & 0 \\
\end{longtable}
\endgroup

\begingroup
\scriptsize
\renewcommand{\arraystretch}{1.08}
\begin{longtable}{@{}c r c p{6.2cm}@{}}
\caption{四问最短时长汇总表}\label{tab:appendix-summary}\\
\toprule
问题 & 最短时长(s) & 最短时长(HH:MM:SS) & 求解状态 \\
\midrule
\endfirsthead
\multicolumn{4}{c}{续表：四问最短时长汇总表}\\
\toprule
问题 & 最短时长(s) & 最短时长(HH:MM:SS) & 求解状态 \\
\midrule
\endhead
\bottomrule
\endfoot
Q1 & 41600 & 11:33:20 & Optimization terminated successfully. (HiGHS Status 7: Optimal) \\
Q2 & 163764 & 45:29:24 & Optimization terminated successfully. (HiGHS Status 7: Optimal) \\
Q3 & 123844 & 34:24:04 & Optimization terminated successfully. (HiGHS Status 7: Optimal) \\
Q4 & 123844 & 34:24:04 & 关键路径下界已达到；最优采购为0 \\
\end{longtable}
\endgroup

